\section{System Model and Problem Formulation}

\subsection{Network Geometry and Channel Model}
We consider the uplink of a Cell-Free Massive MIMO (CF-mMIMO) system comprising $L$ distributed Access Points (APs), indexed by the set $\mathcal{L} = \{1, \dots, L\}$, and $K$ single-antenna user equipments (UEs), indexed by $\mathcal{K} = \{1, \dots, K\}$. Each AP is equipped with $N$ antennas and is connected to a Central Processing Unit (CPU) via a fronthaul link with limited capacity. We assume the network operates in a Time Division Duplex (TDD) mode, where channel reciprocity is exploited for channel estimation.

The wireless channel vector between user $k$ and AP $l$ is denoted by $\mathbf{h}_{lk} \in \mathbb{C}^{N \times 1}$. We adopt a standard block-fading model where the channel remains constant within a coherence block and varies independently between blocks. The channel vector models both large-scale fading (path loss and shadowing) and small-scale fading, given by:
\begin{equation}
    \mathbf{h}_{lk} = \sqrt{\beta_{lk}} \mathbf{g}_{lk},
\end{equation}
where $\beta_{lk} \in \mathbb{R}^+$ represents the large-scale fading coefficient, which is assumed to be known at the CPU and the respective AP due to its slow variation. The vector $\mathbf{g}_{lk} \sim \mathcal{CN}(\mathbf{0}, \mathbf{I}_N)$ represents the small-scale Rayleigh fading component. The spatial correlation between antennas at the same AP is assumed to be negligible.

\subsection{Uplink Transmission and Local Processing}
During the uplink data transmission phase, all $K$ users simultaneously transmit their data symbols to the APs. Let $s_k \in \mathbb{C}$ denote the symbol transmitted by user $k$, satisfying the power constraint $\mathbb{E}[|s_k|^2] = 1$. The transmit power of user $k$ is denoted by $p_k$. The received signal $\mathbf{y}_l \in \mathbb{C}^{N \times 1}$ at AP $l$ is given by:
\begin{equation}
    \mathbf{y}_l = \sum_{k=1}^K \sqrt{p_k} \mathbf{h}_{lk} s_k + \mathbf{n}_l,
\end{equation}
where $\mathbf{n}_l \sim \mathcal{CN}(\mathbf{0}, \sigma^2 \mathbf{I}_N)$ is the additive white Gaussian noise (AWGN) vector at AP $l$ with noise power $\sigma^2$.

In a conventional centralized implementation, APs would forward the raw signals $\mathbf{y}_l$ to the CPU. However, this requires excessive fronthaul bandwidth scaling with $N$. To reduce the dimensionality of the data transmitted over the fronthaul, we assume each AP performs local linear processing to obtain an initial estimate of the user symbols. Specifically, AP $l$ applies a local Linear Minimum Mean Squared Error (LMMSE) combining matrix $\mathbf{W}_l \in \mathbb{C}^{N \times K}$ to the received signal. Assuming perfect local Channel State Information (CSI) is available at the AP, the local estimate $\hat{\mathbf{s}}_l = [\hat{s}_{l1}, \dots, \hat{s}_{lK}]^T \in \mathbb{C}^{K \times 1}$ is obtained as:
\begin{equation}
    \hat{\mathbf{s}}_l = \mathbf{W}_l^H \mathbf{y}_l,
\end{equation}
where the combining matrix is defined as $\mathbf{W}_l = \left( \sum_{k=1}^K p_k \mathbf{h}_{lk} \mathbf{h}_{lk}^H + \sigma^2 \mathbf{I}_N \right)^{-1} \mathbf{H}_l \mathbf{P}^{1/2}$, with $\mathbf{H}_l = [\mathbf{h}_{l1}, \dots, \mathbf{h}_{lK}]$ and $\mathbf{P} = \text{diag}(p_1, \dots, p_K)$. This local estimate $\hat{\mathbf{s}}_l$ serves as the compressed semantic input for the subsequent quantization and neural processing stages.

\subsection{Fronthaul Quantization Constraints}
The local estimates $\hat{\mathbf{s}}_l$ are continuous complex variables that must be quantized before transmission to the CPU. Unlike fixed-rate quantization, we consider a dynamic, task-oriented bit allocation strategy. Let $b_{lk} \in \mathcal{B}$ be the number of bits allocated to quantize the estimate $\hat{s}_{lk}$ for user $k$ at AP $l$, where $\mathcal{B} = \{0, 2, 4\}$ is the set of admissible bit-widths. Here, $b_{lk}=0$ implies that the AP does not forward any information regarding user $k$, effectively performing distributed user selection.

The quantized signal $\tilde{s}_{lk}$ is generated by a quantization function $Q(\cdot)$: 
\begin{equation}
    \tilde{s}_{lk} = Q(\hat{s}_{lk}, b_{lk}).
\end{equation}
The total instantaneous fronthaul rate for AP $l$ is $R_l = \sum_{k=1}^K b_{lk}$. The system is subject to a global average capacity constraint $C_{target}$, reflecting the limited aggregate bandwidth of the fronthaul network.

\subsection{Problem Formulation}
The overarching goal is to recover the transmitted symbol vector $\mathbf{s}$ at the CPU with minimal distortion while satisfying the fronthaul capacity constraints. We formulate this as a joint optimization of the local bit allocation policy $\pi_{\theta}$ (parameterized by $\theta$) and the central reconstruction function $\mathcal{G}_{\phi}$ (parameterized by $\phi$).

The bit allocation policy $\pi_{\theta}$ maps the local observation (the estimate $\hat{\mathbf{s}}_l$ and local channel statistics) to a discrete bit-width $b_{lk}$. The reconstruction function $\mathcal{G}_{\phi}$ aggregates the quantized messages $\{\tilde{\mathbf{s}}_l\}_{l=1}^L$ to produce the final estimate $\hat{\mathbf{s}}_{CPU}$. The optimization problem is formulated as:
\begin{subequations}
\begin{align}
    \min_{\theta, \phi} \quad & \mathbb{E} \left[ \| \mathbf{s} - \hat{\mathbf{s}}_{CPU} \|^2 \right] \\
    \text{s.t.} \quad & \mathbb{E} \left[ \frac{1}{L} \sum_{l=1}^L \sum_{k=1}^K b_{lk} \right] \le C_{target},
\end{align}
\end{subequations}
where the expectation is taken over the channel realizations and noise statistics. This constrained optimization problem balances the trade-off between the detection Mean Squared Error (MSE) and the communication cost. In the proposed method, we relax this problem into an unconstrained Lagrangian formulation to enable end-to-end training via gradient descent.